\chapter{基于 Qwen3 语义先验与门控注意力的统一预测骨架}

本章的主要贡献在于构建了全文后续方法共享的统一建模骨架，并将时间序列预测中的时域建模、基础频域表示与外部冻结语言模型生成的语义先验纳入同一框架中。与仅依赖单一时域编码或单一语义引导的方案相比，本章首先明确了基座模型中“时域分支 + 基础频域分支 + 语义分支 + CMA 对齐 + 共享解码器”的整体数据流，为后续不同频域增强模块的替换提供了统一接口。其次，本文在基座层面引入了基于预测步长的轻量时频融合机制和多头跨模态自适应融合机制，使模型能够在不显著增加复杂度的前提下实现多源信息协同。最后，本章从冻结语言模型主体、仅训练轻量适配模块的角度给出了统一实现路径，使后续第四章和第五章能够在共享训练范式下开展公平对比。  

\section{整体架构设计}

本章提出的统一预测骨架面向多元长程时间序列预测任务，其作用不是直接给出两条最终改进模型，而是建立全文后续方法共享的基座骨架。结合实际实现来看，该骨架由输入归一化模块、时域编码分支、基础频域分支、外部语义嵌入分支、跨模态注意力对齐模块以及共享解码预测头组成。第四章和第五章的两种改进模型均在该骨架上展开，只是分别替换和增强了频域分支及其对应的融合方式。  

设输入多元时间序列为 \(\mathbf{X}\in\mathbb{R}^{B\times L\times C}\)，其中 \(B\) 为 batch 大小，\(L\) 为输入窗口长度，\(C\) 为变量数；外部语言模型生成的变量级语义嵌入记为 \(\mathbf{E}\in\mathbb{R}^{B\times d_{\text{LLM}}\times C}\)，预测步长为 \(H\)。经过 RevIN 归一化后，输入被重排为节点优先形式 \(\mathbf{X}^{\prime}\in\mathbb{R}^{B\times C\times L}\)，以便将每个变量视作一个节点并在统一框架中进行编码和对齐。整体数据流可以概括为以下几个阶段。  

（1）\textbf{时域分支（Time-domain Branch）}：基座模型的时域分支首先利用线性层将历史长度 \(L\) 映射到固定通道维 \(d_{\text{ts}}\)，得到
\[
\mathbf{H}^{\text{time}}_0 = \mathrm{Linear}_{L\rightarrow d_{\text{ts}}}(\mathbf{X}^{\prime}) \in \mathbb{R}^{B\times C\times d_{\text{ts}}}.
\]
随后，将其送入标准 Transformer 编码器以提取变量级时序表示：
\[
\mathbf{H}^{\text{time}} = \mathrm{TFEncoder}(\mathbf{H}^{\text{time}}_0).
\]
这里需要强调的是，第三章对应的基座实现采用的是普通 Transformer 编码器，而不是后续章节中使用的带门控版本。因此，本章的时域分支更适合作为“统一骨架中的基础时域建模模块”来描述，第四章和第五章则在此基础上进一步引入门控机制以增强特征筛选能力。  

（2）\textbf{基础频域分支（Base Frequency Branch）}：在基座模型中，频域信息由一条轻量级频域分支提供。具体而言，模型首先沿时间维对 \(\mathbf{X}^{\prime}\) 进行实数快速傅里叶变换，提取频谱幅值；随后将每个频率 bin 视作一个频率 token，经线性投影映射到通道维后输入频域 Transformer 编码器，并通过注意力池化压缩为变量级频域表示：
\[
\mathbf{H}^{\text{freq}} = \Phi_{\text{fft}}(\mathbf{X}^{\prime}) \in \mathbb{R}^{B\times C\times d_{\text{ts}}}.
\]
因此，第三章中的频域分支本质上是“FFT 幅值建模 + 频率 token 编码 + 频域注意力池化”的基础版本。它的作用是为统一骨架提供一个可被替换的频域接口，而不是最终论文的主要创新点。第四章和第五章正是在这一接口上分别引入小波包多尺度建模与 FreTS 频域前馈网络，以形成两条主要工作路线。  

（3）\textbf{语义分支（Semantic Branch）}：与直接在训练过程中调用大语言模型不同，本文采用外部冻结语言模型生成的变量级语义嵌入作为语义先验输入。记外部嵌入经转置后为
\[
\mathbf{E}' = \mathbf{E}^{\top} \in \mathbb{R}^{B\times C\times d_{\text{LLM}}},
\]
再通过轻量级 \texttt{prompt_encoder} 做结构适配，得到语义特征 \(\mathbf{H}^{\text{sem}}\)。因此，第三章中的语言模型并不直接参与端到端梯度更新，而是以冻结语义先验的形式为后续跨模态对齐提供补充信息。这一表述更符合当前实现，也更便于在论文中清楚区分“外部嵌入生成阶段”和“预测模型训练阶段”。  

（4）\textbf{时频融合机制（Time-Frequency Fusion）}：基座模型中的时域分支与频域分支并非通过复杂的跨注意力结构进行交互，而是采用一种与预测步长相关的轻量门控机制。设时域表示和频域表示在转置后分别为 \(\bar{\mathbf{H}}^{\text{time}}, \bar{\mathbf{H}}^{\text{freq}}\in\mathbb{R}^{B\times d_{\text{ts}}\times C}\)，则模型根据时域编码结果和预测步长 \(H\) 计算门控权重
\[
\mathbf{G}^{\text{hor}} = f_{\text{hor}}(\bar{\mathbf{H}}^{\text{time}}, H)\in\mathbb{R}^{B\times d_{\text{ts}}\times 1},
\]
并完成融合
\[
\mathbf{H}^{\text{tf}} = \mathbf{G}^{\text{hor}} \odot \bar{\mathbf{H}}^{\text{freq}} + (1-\mathbf{G}^{\text{hor}})\odot \bar{\mathbf{H}}^{\text{time}}.
\]
这种设计对应基座模型中的 \texttt{RichHorizonGate}，其作用是根据预测步长动态平衡时域信息与频域信息的贡献。后续两章虽然分别采用更强的融合结构，但其目标仍然与本章一致，即实现时域和频域信息在统一骨架中的协同建模。  

（5）\textbf{跨模态注意力对齐（Cross-modal Alignment）}：在得到时频融合表示 \(\mathbf{H}^{\text{tf}}\) 后，模型将其作为查询序列，将语义特征 \(\mathbf{H}^{\text{sem}}\) 作为键和值，通过多个并行的 CMA 头完成跨模态对齐：
\[
\mathbf{H}^{\text{cma}} = \mathrm{CMA}(\mathbf{H}^{\text{tf}}, \mathbf{H}^{\text{sem}}, \mathbf{H}^{\text{sem}}).
\]
不同头的输出再经 \texttt{AdaptiveDynamicHeadsCMA} 做自适应加权，随后通过残差插值与原始时频表示融合，得到最终的跨模态联合表示 \(\mathbf{H}^{\text{cross}}\)。这一步的核心作用在于利用外部语义先验对变量间关系进行重加权，同时保留数值时序特征本身的稳定结构。  

（6）\textbf{解码与预测头（Decoder and Head）}：最终，联合表示 \(\mathbf{H}^{\text{cross}}\) 被送入 Transformer 解码器，并通过线性层映射到预测步长维度：
\[
\mathbf{Z}^{\text{dec}} = \mathrm{TFDecoder}(\mathbf{H}^{\text{cross}}, \mathbf{H}^{\text{cross}}),\quad
\hat{\mathbf{Y}}^{\text{norm}} = \mathrm{Linear}_{d_{\text{ts}}\rightarrow H}(\mathbf{Z}^{\text{dec}}).
\]
经 RevIN 反归一化后，模型输出最终预测结果 \(\hat{\mathbf{Y}}\in\mathbb{R}^{B\times H\times C}\)。  

综上，本章给出的统一预测骨架可以被视为全文的共享预测基座：它以标准时域编码器、基础 FFT 频域分支和冻结语义先验为起点，通过轻量步长门控与跨模态注意力实现多源信息融合。这样的写法既与基座模型的真实实现保持一致，也为第四章和第五章对频域分支及编码机制的深入增强提供了清晰的参照系。

\section{冻结大模型的适配器设计}

在统一预测骨架中，本文采用“冻结语言模型主体 + 训练预测网络内轻量适配模块”的策略，而不是直接对大语言模型进行端到端微调。更准确地说，变量级文本描述首先由外部冻结语言模型编码为语义嵌入，随后这些嵌入作为输入送入预测网络；在预测网络内部，仅有少量适配器、跨模态对齐模块和下游解码预测模块参与训练。这样的设计可以避免超大规模语言模型反向传播带来的高额计算与显存开销，同时尽可能保留其预训练知识对变量关系建模的辅助作用。  

\subsection{跨模态特征投影与维度对齐}

在跨模态对齐之前，时间序列分支与语义分支处于不同的表示空间。时间序列输入经 RevIN 和维度变换后，形成节点优先表示 \(\mathbf{X}'\in\mathbb{R}^{B\times C\times L}\)；时域分支和频域分支再分别输出统一通道维度的特征，并在基座模型中经过 \texttt{RichHorizonGate} 融合为时频表示 \(\mathbf{H}^{\text{tf}}\)。与此同时，外部冻结语言模型生成变量级语义嵌入 \(\mathbf{E}\in\mathbb{R}^{B\times d_{\text{LLM}}\times C}\)，经转置后得到 \(\mathbf{E}'\in\mathbb{R}^{B\times C\times d_{\text{LLM}}}\)。  

为了在统一空间中进行注意力对齐，本文采用两步式适配策略。第一步，在时间序列侧，通过 \texttt{length_to_feature}、\texttt{freq_token_proj} 等线性映射，将时域输入和频谱 token 投影到统一的通道维 \(d_{\text{ts}}\)，从而保证时域、频域和融合特征在变量维上具有兼容的张量形状。第二步，在语义侧，通过轻量级 \texttt{prompt_encoder} 对 \(\mathbf{E}'\) 进行结构适配，使其能够更好地匹配变量级时序表示的组织方式。  

完成上述对齐后，模型利用 CMA 模块将时频表示作为查询、语义表示作为键和值，即
\[
\mathbf{H}^{\text{cma}} = \mathrm{CMA}\bigl(\mathbf{H}^{\text{tf}},\mathbf{H}^{\text{sem}},\mathbf{H}^{\text{sem}}\bigr).
\]
该过程由多个并行头共同完成，不同头的输出再经 \texttt{AdaptiveDynamicHeadsCMA} 自适应加权融合，形成最终的跨模态对齐表示。最后，模型通过残差插值将其与原始时频表示结合，在保留数值结构稳定性的同时引入语义先验信息。  

\subsection{时序位置编码与语义嵌入设计}

在时间序列分支中，基座模型并未额外引入显式可学习位置编码，时间顺序信息主要通过张量重排和线性映射过程保留下来。具体而言，输入由 \([B,L,C]\) 变换为 \([B,C,L]\)，每个变量节点的历史窗口顺序完整保留，随后通过 \texttt{length_to_feature} 将时间长度映射到固定通道维。也就是说，第三章基座实现中的时间信息更多体现为一种“隐式顺序保持”的编码方式，而不是额外的位置向量注入。  

在语义分支中，语言模型内部的位置编码和语义生成已经在外部嵌入生成阶段完成，预测模型本身并不重新进行 token 级位置建模。本文只对已经得到的变量级语义嵌入使用 \texttt{prompt_encoder} 进行轻量级结构适配，使其在变量维组织方式上更接近时间序列分支的表示。这样可以在不改动语言模型主体的前提下，实现语义特征与数值特征之间的结构对齐。  

这种设计的优点在于职责划分清晰：语言模型负责生成高层语义先验，预测模型负责将该先验与时间序列结构信息进行对齐和融合。相较于在训练阶段直接介入语言模型内部位置编码或嵌入过程，这种做法更符合当前实现，也更有利于控制整体训练复杂度。  

\subsection{适配器参数量与训练效率分析}

与直接微调整个大语言模型不同，本文将可训练参数主要集中在预测网络内部的轻量模块上，具体包括：时间序列侧的 RevIN、\texttt{length_to_feature}、\texttt{freq_token_proj} 等线性映射和编码模块；语义侧的 \texttt{prompt_encoder}；以及跨模态对齐与预测侧的 \texttt{CrossModal}、\texttt{AdaptiveDynamicHeadsCMA}、残差插值参数和 Transformer 解码器。第四章和第五章在此基础上增加的小波包模块或 FreTS 模块，也仍然属于这一类可控规模的扩展参数。  

从训练效率角度看，该设计具有三方面优势。首先，外部语言模型嵌入一旦生成完成，在预测网络训练阶段便无需对语言模型主体存储反向梯度和优化状态，显著降低了显存压力。其次，适配器和时频模块规模相对较小，训练时更容易稳定收敛，也降低了学习率、正则化和微调策略选择的复杂度。最后，由于语言模型主体保持不变，在统一骨架下切换不同频域模块或融合机制时，不需要重新设计整套语义建模流程，从而更便于开展系统性的模型对比和消融分析。  

综合而言，“冻结语言模型主体 + 训练轻量适配器”的策略使本文能够在控制计算成本的前提下利用语言模型的语义先验，并将主要学习负担集中在时序、频域和跨模态对齐模块上。这种设计既符合当前实现，也为后续章节的模块化扩展提供了清晰基础。

\section{门控注意力池化模块}

在统一预测骨架中，时域与频域分支会输出长度不等的序列表示（如每个频段子带对应一段时序、或 FreMLP 后的时域序列），而跨模态对齐与解码器需要固定维度的变量级表示。门控注意力池化模块承担两项职责：其一，在时间或频段维度上通过注意力加权将变长序列压缩为定长向量，即\textbf{注意力池化}；其二，在池化前后或在对多路表示进行融合时引入\textbf{门控机制}，使池化与融合过程具有输入依赖性，从而抑制冗余信息、突出关键时间步或关键频段/头的贡献，与第 2.4 节所述门控注意力的思想一致[门控注意力论文]。  

\textbf{（1）注意力加权池化。} 设某一分支输出的序列表示为 \(\mathbf{H}\in\mathbb{R}^{B\times L\times C}\)（\(B\) 为 batch、\(L\) 为序列长度、\(C\) 为通道维），需在长度维 \(L\) 上聚合为单一向量。采用可学习的注意力池化：先通过小型前馈网络将每个时间步映射为标量得分，再经 Softmax 归一化得到注意力权重，最后对序列做加权求和：  
\[
\mathbf{a} = \mathrm{Softmax}\bigl(f_{\text{attn}}(\mathbf{H})\bigr)\in\mathbb{R}^{B\times L\times 1},\quad
\mathbf{h}_{\text{pooled}} = \sum_{t=1}^{L} a_t\,\mathbf{H}_{:,t,:} \in \mathbb{R}^{B\times C}.
\]
其中 \(f_{\text{attn}}\) 通常为两层 MLP（如 \(\mathrm{Linear}(d\to d/2)\to\mathrm{ReLU}\to\mathrm{Linear}(d/2\to 1)\)）。在小波包版本中，对每个子带分别进行上述注意力池化，得到各频段向量后再通过可学习的频段权重（node_weights）进行加权融合，使不同频段对最终表示的贡献可随数据与任务自适应；在频域 MLP 版本中，对 FreMLP 输出的时域序列直接做一次注意力池化，得到变量级全局表示后再送入后续门控编码器与跨模态模块。  

\textbf{（2）门控机制与 SDPA 输出门控。} 在第四章和第五章的改进模型中，时域或频域分支内部进一步采用带门控的 Transformer 编码层（GatedTransformerEncoderLayer）。对每一层，记归一化后的输入为 \(\tilde{\mathbf{X}}\)，自注意力输出为 \(\mathbf{O}_{\text{attn}}\)，则门控后的输出为  
\[
\mathbf{G} = \sigma\bigl(\mathbf{W}_g\,\tilde{\mathbf{X}}\bigr),\quad
\mathbf{O}_{\text{out}} = \mathbf{O}_{\text{attn}} \odot \mathbf{G},
\]
其中 \(\mathbf{W}_g\) 为可学习线性层，\(\sigma\) 为 Sigmoid 函数，\(\odot\) 为逐元素乘法。该形式与门控注意力文献中“在 SDPA 输出后施加头特定、逐元素的可乘性门控”的设计一致：门控分数由当前层输入决定，具有输入依赖性，并通过 Sigmoid 映射引入稀疏性，从而抑制无关上下文与注意力汇聚现象[门控注意力论文]。因此，在后续改进模型中，池化模块接收到的是经过门控调制、更聚焦于关键时间步的特征；而在第三章基座模型中，虽然编码层尚未完全替换为门控版本，但多头 CMA 的自适应头融合已经体现出“按输入动态选择信息来源”的核心思想。  

\textbf{（3）多头跨模态输出的门控融合。} 在跨模态对齐阶段，框架使用多个独立的 CMA 头（CrossModal），每个头输出变量级特征。为融合多头信息并避免冗余，采用自适应门控融合模块（AdaptiveDynamicHeadsCMA）：将各头输出在通道维拼接后，通过小型 MLP 为每个头生成标量门控分数，经 Softmax 归一化后对多头输出进行加权求和：  
\[
\mathbf{Z} = \mathrm{Concat}\bigl(\mathbf{H}^{\text{cma}}_1,\dots,\mathbf{H}^{\text{cma}}_H\bigr),\quad
\mathbf{g} = \mathrm{Softmax}\bigl(f_{\text{gate}}(\mathbf{Z})\bigr),\quad
\mathbf{H}^{\text{cma}}_{\text{fused}} = \sum_{h=1}^{H} g_h\,\mathbf{H}^{\text{cma}}_h.
\]
该操作等价于在“头”维度上的门控池化：不同样本、不同变量可得到不同的头权重，从而自适应地突出与当前查询更相关的头、抑制冗余头，与门控注意力文献中“头特定门控”带来的收益一致。  

综上，门控注意力池化模块通过\textbf{注意力加权池化}实现变长到定长的压缩，并通过\textbf{多头 CMA 的门控融合}以及后续改进模型中的\textbf{编码层门控}引入输入依赖的选择性，在保证表示能力的同时抑制冗余信息，为跨模态对齐与解码提供更稳定、紧凑的变量级表示，并与第 3.4 节的消融实验相衔接。

\section{基线实验与统一预测骨架有效性验证}

\subsection{实验设置}

为验证第三章所述统一预测骨架的整体有效性，本节首先在与第 2.5 节一致的数据集与划分设置下，构建一组统一的基线实验。所有实验均采用相同的滑动窗口策略，将历史长度 \(L_{\text{in}}\) 的输入序列映射为预测长度 \(L_{\text{pred}}\) 的未来序列，并以 MSE 与 MAE 作为主要评价指标。数据预处理阶段统一使用 RevIN，对训练集、验证集和测试集采用相同的标准化与反标准化流程，以避免由于预处理差异造成结果偏差。  

在模型配置方面，本章采用第三章提出的共享骨架作为默认设置：时域分支使用长度到通道的线性映射以及 Transformer 编码器；语义分支输入由冻结大语言模型离线生成的变量级嵌入组成，并通过轻量级 \texttt{prompt_encoder} 进行结构适配；跨模态对齐模块采用多头 CMA 与自适应头融合；预测端统一采用 Transformer 解码器和线性预测头。为了保证本章比较的可解释性，除专门比较的模块外，其余网络深度、隐藏维度、注意力头数、dropout、优化器、学习率、batch size 与训练轮数均保持一致。  

在具体实现上，建议将第三章实验视为“统一框架有效性验证”，因而需要保持对比对象尽可能简洁：一类为原始基座模型或去掉关键模块的简化模型，另一类为加入冻结语义分支、门控机制和跨模态对齐后的完整统一框架。这样可以保证后续第四章和第五章在切换频域分支时，比较的焦点集中在具体方法本身，而非额外实验设定差异。实验结果可整理为“数据集 × 预测步长 × 模型”的总表，为 3.4.2 和 3.4.3 的分析提供统一依据。

\subsection{Qwen3-1.4B 与 GPT-2 的对比实验}

为验证选择 Qwen3 作为语义骨干的合理性，本节在统一框架下进一步比较不同大语言模型语义先验的作用。对比实验中，时域分支、频域分支、CMA 结构、解码器与训练超参数保持不变，仅替换外部语义嵌入的来源及其对应的输入维度适配层。具体而言，可设置“统一框架 + GPT-2 嵌入”与“统一框架 + Qwen3 嵌入”两组模型，其中两者均采用冻结策略，不参与端到端参数更新。  

该对比的核心目的不是证明更大的语言模型一定在所有场景下都优于较小模型，而是考察不同语义先验与时间序列建模框架的匹配程度。如果 Qwen3 在多数数据集和预测长度上表现更稳定，说明其语义表示在变量关联建模和跨模态对齐方面更适合当前任务；若仅在部分场景下取得优势，则应进一步结合数据集属性分析其原因，例如变量语义信息是否更丰富、序列的跨变量相关性是否更强等。  

在结果写作上，可优先从两个角度展开。第一，从总体指标上比较不同语义骨干下统一框架的 MSE 与 MAE 变化趋势，说明采用 Qwen3 是否能够带来稳定收益。第二，从误差分布或典型数据集角度讨论模型差异，解释语言模型语义先验对后续跨模态注意力重加权的具体影响。这样既能够为第三章选择 Qwen3 提供依据，也能为后续第四、第五章沿用同一语义骨干建立合理前提。

\subsection{门控注意力机制的有效性验证}

为验证第 3.3 节门控注意力池化模块的有效性，本节进一步开展消融实验。考虑到第三章基座模型中已经直接存在的是多头 CMA 输出端的自适应头融合，而编码层内部的门控注意力将在第四章和第五章中体现得更充分，因此本节更适合作为“统一门控思想的预验证”。消融实验可以围绕这两部分展开，例如构造“仅保留多头加权”“将多头输出改为简单平均”“在后续改进模型中进一步比较是否加入编码层门控”等变体。  

通过这种逐步消融，可以较为清楚地回答两个问题：其一，门控是否确实提升了模型预测精度；其二，门控的收益主要来自分支内部的特征筛选，还是来自跨模态阶段的头选择。如果实验结果显示完整模型在 MSE 与 MAE 上均优于对应消融版本，且训练曲线更稳定、波动更小，则可进一步说明门控机制不仅改善了表示质量，也优化了训练过程中的信息流动方式。  

需要强调的是，本节消融的作用在于验证统一框架中的关键设计是否成立，而不是替代第四、第五章对具体频域模块的深入分析。因此，在行文上应将本节的结论收束为三点：统一框架相对原始基座具有整体提升；Qwen3 语义先验为跨模态建模提供了额外收益；门控机制在编码和融合两个层面都具有有效性。基于此，后续两章即可将重点转向两类频域改进路线本身。

\section{本章小结}

本章围绕全文的统一建模骨架，提出了一种面向多元长程时间序列预测的共享改进框架。该框架以 RevIN 预处理为起点，在时域分支中提取变量级动态特征，在语义分支中利用冻结大语言模型生成的外部嵌入作为先验，并通过跨模态注意力模块实现时序特征与语义特征的对齐。与此同时，本文在框架中引入门控注意力池化机制，使模型能够在编码、池化和跨头融合三个层面抑制冗余信息、突出关键表示。  

本章的意义在于为后续两条主要技术路线提供统一出发点。第四章将在该骨架下将频域分支具体化为小波包分解增强模块，重点讨论多尺度子带对大误差抑制与 MSE 优化的作用；第五章则将频域分支替换为可学习的频域前馈网络，重点分析复数域频谱建模、稀疏化与步长感知门控融合对平均误差和长期趋势预测的改善效果。
